% GENERATED FILE. Edit canonical lessons, then run npm run generate:books.
% canonical-source-hash: fnv1a64:b717cb0537507423
% canonical-lessons: SA-C432-vyayah, SA-C432-apavyayah, SA-C432-vyayati, SA-C432-sancinoti, SA-C432-sancayah

\chapter{Spending and Saving}
\label{ch:sa-spending-and-saving}
\hlchaptermodality{432}

\begin{chapteropening}
\textbf{By the end of this chapter:} \emph{I can say what I spend, earn and save on a journey.}

Complete the last lesson of chapter 432: I can say what I spend, earn and save on a journey.
\end{chapteropening}

\section[\texorpdfstring{\sk{व्ययः} (vyayaḥ)}{व्ययः (vyayaḥ)}]{\sk{व्ययः} (vyayaḥ) — an expense}

\label{lesson:SA-C432-vyayah}



\begin{quote}
Before the new one: say the Sanskrit for mixes, then the Sanskrit for a touchstone.
\end{quote}

\subsection*{You'll want to know: \sk{व्ययः}}
\textbf{\sk{व्ययः}} — \emph{vyayaḥ} — \textquotedblleft{}an expense\textquotedblright{}.

\begin{quote}
\sk{तीर्थयात्रायां} \sk{महान्} \sk{व्ययः} \sk{भवति।}
\end{quote}

\begin{quote}
\emph{tīrthayātrāyāṃ mahān vyayaḥ bhavati.}
\end{quote}

\begin{quote}
A pilgrimage brings a big expense.
\end{quote}

A masculine noun.

\begin{cousinweb}
\textbf{\sk{व्ययः}} (\emph{vyayaḥ}) is built from \emph{vi}, \textquotedblleft{}away\textquotedblright{}, and the root √i, \textquotedblleft{}to go\textquotedblright{}: a going away of money; MW gives \emph{vyaya} as spending, expense (MW).
\end{cousinweb}

\begin{grammarlens}[title={What you've built}]
One more word for getting around.
\end{grammarlens}

\subsection*{Your turn}
\begin{itemize}
\raggedright
  \item \emph{Say it:} \emph{vyayaḥ}
  \item \emph{Say it:} \emph{vyayaḥ}, once more
  \item \emph{Say it:} \emph{tīrthayātrāyāṃ mahān vyayaḥ bhavati.} — A pilgrimage brings a big expense.
  \item \emph{Recall:} say the Sanskrit for mixes, then the Sanskrit for a touchstone, then say \emph{vyayaḥ} again
\end{itemize}

\begin{tcolorbox}[breakable,colback=teal!4,colframe=teal!35!black,title={Before you move on}]
What does \sk{व्ययः} mean? (\textquotedblleft{}An expense\textquotedblright{}.) Say it once more.
\end{tcolorbox}
\section[\texorpdfstring{\sk{अपव्ययः} (apavyayaḥ)}{अपव्ययः (apavyayaḥ)}]{\sk{अपव्ययः} (apavyayaḥ) — wasteful spending}

\label{lesson:SA-C432-apavyayah}



\begin{quote}
Before the new one: say the Sanskrit for a touchstone, then the Sanskrit for an expense.
\end{quote}

\subsection*{You'll want to know: \sk{अपव्ययः}}
\textbf{\sk{अपव्ययः}} — \emph{apavyayaḥ} — \textquotedblleft{}wasteful spending\textquotedblright{}.

\begin{quote}
\sk{अपव्ययः} \sk{उचितः} \sk{न} \sk{अस्ति।}
\end{quote}

\begin{quote}
\emph{apavyayaḥ ucitaḥ na asti.}
\end{quote}

\begin{quote}
Wasteful spending is not right.
\end{quote}

A masculine noun.

\begin{cousinweb}
\textbf{\sk{अपव्ययः}} (\emph{apavyayaḥ}) is built from \emph{apa}, \textquotedblleft{}away, wrongly\textquotedblright{}, and \emph{vyaya}, \textquotedblleft{}expense\textquotedblright{}; MW gives \emph{apavyaya} as prodigality (MW).
\end{cousinweb}

\begin{grammarlens}[title={What you've built}]
One more word for getting around.
\end{grammarlens}

\subsection*{Your turn}
\begin{itemize}
\raggedright
  \item \emph{Say it:} \emph{apavyayaḥ}
  \item \emph{Say it:} \emph{apavyayaḥ}, once more
  \item \emph{Say it:} \emph{apavyayaḥ ucitaḥ na asti.} — Wasteful spending is not right.
  \item \emph{Recall:} say the Sanskrit for a touchstone, then the Sanskrit for an expense, then say \emph{apavyayaḥ} again
\end{itemize}

\begin{tcolorbox}[breakable,colback=teal!4,colframe=teal!35!black,title={Before you move on}]
What does \sk{अपव्ययः} mean? (\textquotedblleft{}Wasteful spending\textquotedblright{}.) Say it once more.
\end{tcolorbox}
\section[\texorpdfstring{\sk{व्ययति} (vyayati)}{व्ययति (vyayati)}]{\sk{व्ययति} (vyayati) — spends; lays out money}

\label{lesson:SA-C432-vyayati}



\begin{quote}
Before the new one: say the Sanskrit for an expense, then the Sanskrit for wasteful spending.
\end{quote}

\subsection*{You'll want to know: \sk{व्ययति}}
\textbf{\sk{व्ययति}} — \emph{vyayati} — \textquotedblleft{}spends; lays out money\textquotedblright{}.

\begin{quote}
\sk{सः} \sk{यात्रायां} \sk{बहु} \sk{धनं} \sk{व्ययति।}
\end{quote}

\begin{quote}
\emph{saḥ yātrāyāṃ bahu dhanaṃ vyayati.}
\end{quote}

\begin{quote}
He spends a lot of money on the journey.
\end{quote}

Past (he …): \sk{सः} \sk{यात्रायां} \sk{धनं} \sk{व्ययितवान्}, \emph{saḥ yātrāyāṃ dhanaṃ vyayitavān}.

\begin{cousinweb}
\textbf{\sk{व्ययति}} (\emph{vyayati}) comes from the root √vyay, \textquotedblleft{}to spend\textquotedblright{}, which MW takes rather as a verb made from \emph{vyaya}, \textquotedblleft{}an expense\textquotedblright{}; MW gives \emph{vyayati} as to expend or spend, and KRM gives its participle \emph{vyayita} (MW).
\end{cousinweb}

\begin{grammarlens}[title={What you've built}]
One more word for getting around.
\end{grammarlens}

\subsection*{Your turn}
\begin{itemize}
\raggedright
  \item \emph{Say it:} \emph{vyayati}
  \item \emph{Say it:} \emph{vyayati}, once more
  \item \emph{Say it:} \emph{saḥ yātrāyāṃ bahu dhanaṃ vyayati.} — He spends a lot of money on the journey.
  \item \emph{Recall:} say the Sanskrit for an expense, then the Sanskrit for wasteful spending, then say \emph{vyayati} again
\end{itemize}

\begin{tcolorbox}[breakable,colback=teal!4,colframe=teal!35!black,title={Before you move on}]
What does \sk{व्ययति} mean? (\textquotedblleft{}Spends; lays out money\textquotedblright{}.) Say it once more.
\end{tcolorbox}
\section[\texorpdfstring{\sk{सञ्चिनोति} (sañcinoti)}{सञ्चिनोति (sañcinoti)}]{\sk{सञ्चिनोति} (sañcinoti) — saves up; gathers}

\label{lesson:SA-C432-sancinoti}



\begin{quote}
Before the new one: say the Sanskrit for wasteful spending, then the Sanskrit for spends.
\end{quote}

\subsection*{You'll want to know: \sk{सञ्चिनोति}}
\textbf{\sk{सञ्चिनोति}} — \emph{sañcinoti} — \textquotedblleft{}saves up; gathers\textquotedblright{}.

\begin{quote}
\sk{सः} \sk{प्रतिदिनं} \sk{किञ्चित्} \sk{धनं} \sk{सञ्चिनोति।}
\end{quote}

\begin{quote}
\emph{saḥ pratidinaṃ kiñcit dhanaṃ sañcinoti.}
\end{quote}

\begin{quote}
Every day he saves up a little money.
\end{quote}

Past (he …): \sk{सः} \sk{धनं} \sk{सञ्चितवान्}, \emph{saḥ dhanaṃ sañcitavān}.

\begin{cousinweb}
\textbf{\sk{सञ्चिनोति}} (\emph{sañcinoti}) is built from \emph{sam}, \textquotedblleft{}together\textquotedblright{}, and the root √ci, \textquotedblleft{}to gather, to heap\textquotedblright{}; MW gives \emph{saṃ-ci} as to heap up, accumulate, with the participle \emph{saṃcita}, \textquotedblleft{}collected\textquotedblright{} (MW).
\end{cousinweb}

\begin{grammarlens}[title={What you've built}]
One more word for getting around.
\end{grammarlens}

\subsection*{Your turn}
\begin{itemize}
\raggedright
  \item \emph{Say it:} \emph{sañcinoti}
  \item \emph{Say it:} \emph{sañcinoti}, once more
  \item \emph{Say it:} \emph{saḥ pratidinaṃ kiñcit dhanaṃ sañcinoti.} — Every day he saves up a little money.
  \item \emph{Recall:} say the Sanskrit for wasteful spending, then the Sanskrit for spends, then say \emph{sañcinoti} again
\end{itemize}

\begin{tcolorbox}[breakable,colback=teal!4,colframe=teal!35!black,title={Before you move on}]
What does \sk{सञ्चिनोति} mean? (\textquotedblleft{}Saves up; gathers\textquotedblright{}.) Say it once more.
\end{tcolorbox}
\section[\texorpdfstring{\sk{सञ्चयः} (sañcayaḥ)}{सञ्चयः (sañcayaḥ)}]{\sk{सञ्चयः} (sañcayaḥ) — savings; a store}

\label{lesson:SA-C432-sancayah}



\begin{quote}
Before the new one: say the Sanskrit for spends, then the Sanskrit for saves up.
\end{quote}

\subsection*{You'll want to know: \sk{सञ्चयः}}
\textbf{\sk{सञ्चयः}} — \emph{sañcayaḥ} — \textquotedblleft{}savings; a store\textquotedblright{}.

\begin{quote}
\sk{मम} \sk{सञ्चयः} \sk{वित्तकोषे} \sk{अस्ति।}
\end{quote}

\begin{quote}
\emph{mama sañcayaḥ vittakoṣe asti.}
\end{quote}

\begin{quote}
My savings are in the bank.
\end{quote}

A masculine noun.

\begin{cousinweb}
\textbf{\sk{सञ्चयः}} (\emph{sañcayaḥ}) is built from \emph{sam}, \textquotedblleft{}together\textquotedblright{}, and the root √ci, \textquotedblleft{}to gather\textquotedblright{}; MW gives \emph{saṃcaya} as collection, accumulation, a store (MW).
\end{cousinweb}

\begin{grammarlens}[title={What you've built}]
One more word for getting around.
\end{grammarlens}

\subsection*{Your turn}
\begin{itemize}
\raggedright
  \item \emph{Say it:} \emph{sañcayaḥ}
  \item \emph{Say it:} \emph{sañcayaḥ}, once more
  \item \emph{Say it:} \emph{mama sañcayaḥ vittakoṣe asti.} — My savings are in the bank.
  \item \emph{Recall:} say the Sanskrit for spends, then the Sanskrit for saves up, then say \emph{sañcayaḥ} again
\end{itemize}

\begin{tcolorbox}[breakable,colback=teal!4,colframe=teal!35!black,title={Before you move on}]
What does \sk{सञ्चयः} mean? (\textquotedblleft{}Savings; a store\textquotedblright{}.) Say it once more.
\end{tcolorbox}
